% Preamble for master-thesis PZ build (trimmed from thesis/upstream-template).

\usepackage{polyglossia}

\usepackage{fontspec}
\defaultfontfeatures{Ligatures=TeX}

\newcommand{\setmainthesisfont}[1]{%
  \setmainfont{#1}[Ligatures=TeX]%
  \newfontfamily\russianfont{#1}[Script=Cyrillic,Ligatures=TeX]%
  \newfontfamily\cyrillicfont{#1}[Script=Cyrillic,Ligatures=TeX]%
}
\newcommand{\setsansthesisfont}[1]{%
  \setsansfont{#1}[Ligatures=TeX,Scale=MatchLowercase]%
  \newfontfamily\cyrillicfontsf{#1}[Script=Cyrillic,Ligatures=TeX,Scale=MatchLowercase]%
}
\newcommand{\setmonothesisfont}[1]{%
  \setmonofont{#1}[Scale=MatchLowercase]%
  \newfontfamily\cyrillicfonttt{#1}[Scale=MatchLowercase]%
}

\IfFontExistsTF{Times New Roman}{%
  \setmainthesisfont{Times New Roman}%
  \setmonothesisfont{Times New Roman}%
}{%
  \setmainthesisfont{Tinos}%
  \setmonothesisfont{Tinos}%
}
\IfFontExistsTF{Arial}{%
  \setsansthesisfont{Arial}%
}{%
  \setsansthesisfont{Arimo}%
}

\setmainlanguage[numerals=cyrillic]{russian}
\setotherlanguages{english}
\usepackage{csquotes}

\usepackage{xunicode} % some extra unicode support
%\usepackage[utf8x]{inputenc}
\usepackage{xltxtra} % \XeLaTeX macro

%\setromanfont{Charis SIL}
%\setsansfont{Liberation Sans}
%\setmonofont{PT Mono}
%\setmainfont{Liberation Serif} % this allows to use sans-serif as default font

%нумерация справа и колонтитулы справа вверху
\usepackage{fancyhdr}
\usepackage[a4paper,left=30mm,right=15mm,top=20mm,bottom=20mm,bindingoffset=0cm]{geometry}%
\setlength{\parindent}{1.25cm}
\setlength{\parskip}{0pt}

\usepackage{indentfirst}


\usepackage[protrusion=true]{microtype}
\emergencystretch 3em
\hfuzz=0.5pt

\usepackage{mathtools}
\usepackage{amsfonts}
\usepackage{amssymb}
\usepackage{amsmath}
\usepackage{amsthm}

\usepackage{calc}
\usepackage{ifthen}
\usepackage{graphicx}
\usepackage{tabularx}
\usepackage{array}
\usepackage{ragged2e}
\AtBeginDocument{\normalfont\normalsize\justifying}
\newcolumntype{P}[1]{>{\RaggedRight\hspace{0pt}\arraybackslash}p{#1}}
\usepackage{longtable}
\usepackage{needspace}
\usepackage{multirow}
\usepackage{float}
\usepackage{placeins}
\makeatletter %попытка избавиться от нежелательных разрывов строк, вроде не очень работает
\def\nobreakhline{%
  \noalign{\ifnum0=`}\fi
    \penalty\@M
    \futurelet\@let@token\LT@@nobreakhline}
\def\LT@@nobreakhline{%
  \ifx\@let@token\hline
    \global\let\@gtempa\@gobble
    \gdef\LT@sep{\penalty\@M\vskip\doublerulesep}% <-- change here
  \else
    \global\let\@gtempa\@empty
    \gdef\LT@sep{\penalty\@M\vskip-\arrayrulewidth}% <-- change here
  \fi
  \ifnum0=`{\fi}%
  \multispan\LT@cols
     \unskip\leaders\hrule\@height\arrayrulewidth\hfill\cr
  \noalign{\LT@sep}%
  \multispan\LT@cols
     \unskip\leaders\hrule\@height\arrayrulewidth\hfill\cr
  \noalign{\penalty\@M}%
  \@gtempa}
\makeatother

\usepackage[unicode=true,hidelinks]{hyperref}
\usepackage{xurl}
\urlstyle{same}
\Urlmuskip=0mu
\usepackage{color}
\usepackage{pgf}

\usepackage{titling}

% Настройка списков (без лишних вертикальных отступов)
\usepackage{paralist}
\setdefaultenum{1.}{1.}{1.}{1.}
\setdefaultitem{--}{}{}{}
%\setlength\itemsep{-1em}
\let\itemize\compactitem
\let\enditemize\endcompactitem
\let\enumerate\compactenum
\let\endenumerate\endcompactenum
\let\description\compactdesc
\let\enddescription\endcompactdesc
\pltopsep=\smallskipamount
\plitemsep=0pt
\plparsep=0pt
% Команда для отмены разрыва страниц перед списками
\makeatletter 
\newcommand\mynobreakpar{\par\nobreak\@afterheading} 
\makeatother
%%%%%%

\usepackage[singlelinecheck=true]{caption}
\DeclareCaptionLabelSeparator{thesisdash}{~---~}
\captionsetup{labelsep=thesisdash}
\usepackage{subcaption}
\captionsetup[table]{justification=justified,singlelinecheck=false,name=Таблица,position=top}
\captionsetup[figure]{justification=justified,name=Рисунок,singlelinecheck=on,font=onehalfspacing}

\usepackage{titlesec}
%\titleformat{\sectionname}[runin]{начертание}{счетчик, например \thesubparagraph}{0pt}{}{}

\titleformat{\chapter}[block]{\centering\normalfont\fontsize{14pt}{17pt}\selectfont\bfseries}{\thechapter.}{1ex}{}{}
\titlespacing{\chapter}{0pt}{0em}{0pt}

\titleformat{\section}[block]{\normalfont\fontsize{13pt}{16pt}\selectfont\bfseries}{\thesection}{1ex}{}{}
\titlespacing{\section}{0pt}{0em}{0pt}

\titleformat{\subsection}[block]{\normalfont\fontsize{12pt}{15pt}\selectfont\bfseries}{\thesubsection}{1ex}{}{}
\titlespacing{\subsection}{0pt}{0em}{0pt}

	% paragraph и subparagraph -- в тексте, без отступов
% \titlespacing{command}{left spacing}{before spacing}{after spacing}[right]

\titleformat{\paragraph}[runin]{\normalfont\normalsize\bfseries}{\theparagraph}{0pt}{}{}
\titlespacing{\paragraph}{0pt}{0em}{0ex} %не добавлять дополнительных отступов
% \titlespacing{\paragraph}{\parindent}{0em}{10ex}[1em] %добавить отступы

\titleformat{\subparagraph}[runin]{\normalfont\normalsize\bfseries}{\thesubparagraph}{0pt}{}{}
\titlespacing{\subparagraph}{0pt}{0em}{0ex}


% Своё название для Cписка литературы
\usepackage[title, titletoc]{appendix}
\addto\captionsrussian{% Replace "english" with the language you use
	\renewcommand{\contentsname}%
	{Содержание}%
}

%\renewcommand{\appendixname}{Приложение}% Change "chapter name" for Appendix chapters
%\renewcommand{\cftchapdotsep}{\cftdotsep}

\pagestyle{fancy}
\fancyhf{}
\fancyfoot[C]{{\normalfont\normalsize\rmfamily\thepage}}
\renewcommand{\headrulewidth}{0pt}
\renewcommand{\footrulewidth}{0pt}
\fancypagestyle{plain}{%
  \fancyhf{}%
  \fancyfoot[C]{{\normalfont\normalsize\rmfamily\thepage}}%
  \renewcommand{\headrulewidth}{0pt}%
  \renewcommand{\footrulewidth}{0pt}%
}
\renewcommand{\baselinestretch}{1.5}
\newcommand{\headertext}[1]{\fancyhead[R]{\tiny{#1}}}


% поддержка нумерованных списков с нумерацией в виде русских букв
\usepackage{enumitem}
\makeatletter
\AddEnumerateCounter{\asbuk}{\russian@asbuk@alph}{а} % section 8.1
\renewcommand{\thesubfigure}{\asbuk{subfigure}}
\makeatother

%% Список литературы

\usepackage[
  style=gost-numeric,
  sorting=none,
  language=auto,
  autolang=other
]{biblatex}
\addbibresource{chapters/biblio.bib}
\renewcommand\multicitedelim{,\space}

\usepackage{tikz}
\usetikzlibrary{shapes.geometric,arrows.meta,positioning}
\usepackage{hhline}

%\frenchspacing %% изменение расстояние до и после точек в ряде случаев

\renewcommand{\theenumi}{\arabic{enumi}}
\renewcommand{\theenumii}{\arabic{enumii}}
\renewcommand{\theenumiii}{\arabic{enumiii}}
\renewcommand{\theenumiv}{\arabic{enumiv}}

\renewcommand{\labelenumi}{\theenumi.}
\renewcommand{\labelenumii}{\theenumi.\theenumii.}
\renewcommand{\labelenumiii}{\theenumi.\theenumii.\theenumiii.}
\renewcommand{\labelenumiv}{\theenumi.\theenumii.\theenumiii.\theenumiv.}

%\newenvironment{annotation}{\textbf{Аннотация.} \textit}{}
\theoremstyle{plain}
\newtheorem*{annotation}{Аннотация}

\headertext{}

\usepackage{listings}

\renewcommand{\lstlistingname}{Листинг}

\lstset{
  basicstyle=\ttfamily\small,
  tabsize=2,
  showstringspaces=false,
  columns=fullflexible,
  keepspaces=true,
  numbers=left,
  stepnumber=1,
  numberstyle=\tiny\color{gray},
  breaklines=true,
  breakatwhitespace=false,
  framesep=6pt,
  aboveskip=0.5em,
  belowskip=0.5em,
  lineskip=0pt,
  abovecaptionskip=1em,
  captionpos=t,
  extendedchars=true,
  literate={Ö}{{\"O}}1
  {Ä}{{\"A}}1
  {Ü}{{\"U}}1
  {ß}{{\ss}}1
  {ü}{{\"u}}1
  {ä}{{\"a}}1
  {ö}{{\"o}}1
  {~}{{\textasciitilde}}1
  {…}{\ldots}1
  {–}{-}1
  {\ }{ }1
}

\newcounter{totalfigures}
\newcounter{totaltables}
\newcounter{totallistings}

%%% Local Variables:
%%% TeX-engine: xetex
%%% End:
